\documentclass{article}
\usepackage{spconf,amsmath,graphicx}
\usepackage[utf8]{inputenc}
\usepackage[T1]{fontenc}
\usepackage[brazil]{babel}
\usepackage{enumitem}
\usepackage{csquotes}
\usepackage{hyperref}
\hypersetup{
    colorlinks=true,
    linkcolor=black,
    urlcolor=blue,
    citecolor=black
}

% Title.
% ------
\title{Projeto Integrador I --- Relatório 1}
%
% Single address.
% ---------------
\name{Gabriel Augusto Camolezi}
\address{Universidade Federal de Santa Maria (UFSM) \\
         Campus Cachoeira do Sul \\
         Projeto Integrador I --- 25 de setembro de 2026}

\begin{document}
\maketitle
%
\begin{abstract}
Este relatório apresenta o projeto de análise e previsão dos agendamentos e dos não comparecimentos às refeições nos Restaurantes Universitários (RUs) da Universidade Federal de Santa Maria (UFSM). A partir de dados históricos, pretende-se identificar padrões de comportamento dos usuários e investigar métodos de previsão que possam auxiliar no planejamento da produção de refeições, reduzindo o desperdício de alimentos. São descritos o contexto, a motivação, os objetivos geral e específicos do projeto, além do cronograma de atividades previsto para as próximas quatro semanas.
\end{abstract}
%
\begin{keywords}
Restaurantes universitários, previsão de demanda, não comparecimento, análise de dados, UFSM
\end{keywords}
%
\section{Introdução}
\label{sec:intro}

\subsection{Contexto geral}
\label{ssec:contexto}

Os Restaurantes Universitários (RUs) da Universidade Federal de Santa Maria (UFSM) atendem uma quantidade significativa de usuários por meio do agendamento de refeições. Nesse contexto, existe a necessidade de compreender o comportamento dos agendamentos e dos não comparecimentos às refeições previamente agendadas.

O projeto está inserido no contexto de análise de dados e modelagem preditiva, buscando utilizar dados históricos relacionados aos RUs para identificar padrões e construir métodos capazes de estimar a quantidade de usuários que realizarão agendamentos e a quantidade de usuários que não comparecerão às refeições agendadas.

A pergunta de pesquisa que orienta o projeto é:

\begin{quote}
\itshape
É possível prever a quantidade de usuários que agendarão refeições e a quantidade de usuários que não comparecerão aos Restaurantes Universitários da UFSM?
\end{quote}

\subsection{Motivo da relevância do projeto}
\label{ssec:relevancia}

A relevância do projeto está relacionada à possibilidade de utilizar a análise preditiva para auxiliar no planejamento da quantidade de refeições que devem ser preparadas nos Restaurantes Universitários da UFSM. Caso as previsões de agendamentos e de não comparecimentos apresentem resultados suficientemente adequados, essas informações poderão ser consideradas no planejamento da produção de refeições.

Dessa forma, uma estimativa mais precisa da quantidade de usuários que efetivamente comparecerão aos Restaurantes Universitários poderá contribuir para reduzir a produção excessiva de refeições e, consequentemente, evitar o desperdício de alimentos na Universidade.

\subsection{O que se pretende alcançar}
\label{ssec:objetivo-geral-intro}

Pretende-se verificar, por meio da análise dos dados disponíveis, se é possível realizar previsões sobre a quantidade de usuários que agendarão refeições e sobre a quantidade de usuários que não comparecerão aos RUs da UFSM.

Para isso, o projeto deverá explorar os dados, identificar padrões relevantes, analisar possíveis relações entre diferentes variáveis e testar métodos de previsão que possam ser aplicados ao problema.

\section{Motivação}
\label{sec:motivacao}

A motivação do projeto está relacionada à necessidade de compreender melhor o comportamento dos agendamentos e dos não comparecimentos nos Restaurantes Universitários da UFSM. A disponibilidade de dados históricos possibilita investigar esse comportamento, buscando transformar os registros existentes em informações que possam auxiliar na compreensão da demanda.

Além disso, o problema apresenta uma oportunidade de aplicação prática dos conhecimentos de Ciência de Dados e Inteligência Artificial, envolvendo etapas de obtenção e preparação dos dados, análise exploratória, visualização, identificação de padrões e desenvolvimento de modelos preditivos.

\section{Objetivo Geral}
\label{sec:objgeral}

Investigar a possibilidade de prever a quantidade de usuários que agendarão refeições e a quantidade de usuários que não comparecerão aos Restaurantes Universitários da UFSM, utilizando dados históricos e técnicas de análise e modelagem preditiva.

\section{Objetivos Específicos}
\label{sec:objespec}

\begin{enumerate}[label=\arabic*.]
    \item \textbf{Levantar e organizar os dados} relacionados aos agendamentos, refeições e não comparecimentos dos Restaurantes Universitários da UFSM.

    \item \textbf{Realizar uma análise exploratória dos dados}, utilizando estatísticas descritivas e visualizações para compreender o comportamento dos agendamentos e dos não comparecimentos.

    \item \textbf{Investigar padrões e possíveis relações} entre os agendamentos, os não comparecimentos e outras variáveis disponíveis na base de dados.

    \item \textbf{Avaliar a influência de fatores relevantes}, incluindo variáveis temporais e, quando houver dados adequados, informações relacionadas às condições meteorológicas e às avaliações dos Restaurantes Universitários.

    \item \textbf{Selecionar e aplicar métodos de previsão} adequados ao problema, buscando estimar a quantidade de agendamentos e de não comparecimentos.

    \item \textbf{Avaliar os resultados dos métodos utilizados}, verificando seu desempenho e suas limitações para o contexto dos Restaurantes Universitários da UFSM.

    \item \textbf{Documentar os procedimentos e resultados obtidos}, apresentando as conclusões da investigação e as perspectivas de continuidade do projeto.
\end{enumerate}

\section{Cronograma --- Próximas Quatro Semanas}
\label{sec:cronograma}

\subsection{Até 02/10 --- Visualização dos dados}
\label{ssec:etapa1}

\begin{itemize}
    \item Reproduzir em Python, utilizando \texttt{Matplotlib} e \texttt{Jupyter Notebook}, os gráficos correspondentes às visualizações já desenvolvidas em PostgreSQL a partir das bases de dados de avaliações e acessos dos Restaurantes Universitários.
    \item Comparar os resultados obtidos em Python com os resultados das consultas e visualizações realizadas no PostgreSQL.
    \item Organizar os gráficos produzidos para utilização nos próximos relatórios e análises do projeto.
\end{itemize}

\subsection{Até 09/10 --- Geração de dados sintéticos}
\label{ssec:etapa2}

\begin{itemize}
    \item Analisar os dados reais de avaliações dos Restaurantes Universitários referentes a 2022, identificando os atributos relevantes para a geração dos dados sintéticos.
    \item Calcular medidas estatísticas dos atributos selecionados, como média e desvio-padrão, e identificar as respectivas distribuições de probabilidade.
    \item Definir, a partir dessas características estatísticas, um método de amostragem para gerar dados sintéticos referentes ao período de 2023 a 2026.
    \item Gerar e validar os dados sintéticos, verificando se suas principais características estatísticas são compatíveis com as observadas na base real de 2022.
    \item Documentar o processo de geração e identificar claramente os dados produzidos como sintéticos.
\end{itemize}

\subsection{Até 16/10 --- Segundo relatório do Projeto Integrador I}
\label{ssec:etapa3}

\begin{itemize}
    \item Consolidar os resultados obtidos nas etapas anteriores.
    \item Redigir o segundo relatório do Projeto Integrador I.
    \item Revisar a metodologia, os resultados e as análises realizadas.
    \item Finalizar e entregar o relatório até a data estabelecida.
\end{itemize}

\subsection{Até 23/10 --- Investigar relação entre chuva e comportamento dos usuários}
\label{ssec:etapa4}

\begin{itemize}
    \item Integrar os dados meteorológicos do Instituto Nacional de Meteorologia (INMET) aos dados dos Restaurantes Universitários.
    \item Investigar a existência de correlação entre a ocorrência e/ou intensidade de chuva e variáveis relacionadas ao funcionamento dos RUs, como quantidade de agendamentos, acessos e não comparecimentos.
    \item Definir a escala temporal adequada para a análise, considerando especialmente a relação entre data, horário da refeição e ocorrência de chuva.
    \item Calcular e interpretar as correlações encontradas, distinguindo associação estatística de relação causal.
    \item Produzir visualizações que auxiliem na interpretação dos resultados.
\end{itemize}

\end{document}
